\section{Prime-to-\texorpdfstring{$p$}{p} roots from stable reduction}
\label{sec:tame-roots}

We first treat the extensions whose horizontal ramification indices divide
the fixed prime $\ell\ne p$.  The patching construction reduces this task to
extracting $\ell$-th roots of coordinates on collars around the branch
points.  We shall find these roots in finite \'etale covers of $X_\Omega$.
There are two steps: extract roots on stable node annuli, and use resolution
of nonsingularities to reach such annuli from arbitrary small disks.

\subsection{Units and roots on node annuli}

We use the following elementary analytic fact.  The exponent in its
factorization is the slope of the additive valuation of the unit along the
annulus skeleton.
The unit factorization is the standard annular factorization; compare
\cite[Proposition~2.2(1)]{BakerPayneRabinoff}.  We give the argument on
the entire open annulus, since this is the domain required for gluing.

\begin{lemma}[Units on disks and annuli]\label{lem:annular-unit}
An analytic unit on an open disk over $\Omega$ has constant absolute value
in $|\Omega^\times|$.  If $A=\{r<|S|<R\}$ is an open annulus and
$u\in\mathcal O(A)^\times$, then
\[
 u=cS^d(1+h),\qquad c\in\Omega^\times,\quad d\in\mathbb Z,\quad
 h\in\mathcal O(A),\quad |h(x)|<1\quad(x\in A),
\]
where $cS^d$ is a term of the Laurent expansion of $u$.
The function $1+h$ has an analytic $\ell$-th root on $A$.
Consequently, on any analytic cover of $A$ on which $S$ has an
$\ell$-th root, so does $u$.
\end{lemma}

\begin{proof}
On a circle $|S|=\rho$ of radius $\rho\in|\Omega^\times|$, the Laurent
Gauss norm is multiplicative.  Normalize $u$ and its inverse to have norm
one.  Their reductions are inverse units in
$\widetilde\Omega[s,s^{-1}]$, so the reduction of $u$ is a monomial.
Thus exactly one Laurent term dominates at that radius.  Its exponent
cannot change inside the radius interval.  Indeed, convergence on compact
subannuli reduces any proposed change to finitely many competing terms;
a tie between two terms occurs at a radius in $|\Omega^\times|$, since
the additive value group is $\mathbb Q$.  Such a tie contradicts the
preceding reduction argument.  A single term $cS^d$ therefore dominates
throughout $A$, giving the asserted factorization.  The same argument on
closed subdisks uses that the only units in $\widetilde\Omega[s]$ are
constants and proves the disk assertion.  Its constant norm belongs to
$|\Omega^\times|$ by evaluation at a rational point of the disk.

Since $\ell$ is a unit in the valuation ring, the binomial series for
$(1+h)^{1/\ell}$ converges on every compact subannulus; the resulting
functions agree on overlaps.  A root of $c$ lies in $\Omega$, and a root
of $S^d$ follows from a root of $S$, also when $d<0$.
\end{proof}

Recall the data attached to a split stable model of a proper curve.
Its skeleton is a finite metric graph.  Each vertex carries the smooth
normalization of the corresponding special-fibre component, with one
marked point for each incident edge end.  A loop has two distinct marks
on the same normalization.  The tube of a node is an open annulus,
and the tube of a smooth special-fibre point is an open disk.
The graph with these marked residue curves is its \emph{metrized complex}.
Stability means
\[
 2g(v)-2+\operatorname{val}(v)>0
\]
at each vertex $v$, where $g(v)$ is the genus of its residue curve.

\begin{lemma}[Roots on stable node annuli]\label{lem:stable-edge-root}
Let $C/\Omega$ be a smooth proper connected curve of genus at least two,
equipped with its stable skeleton.  Let $A$ be a node annulus and $S$
an annular coordinate on $A$.  There is a connected finite \'etale cover
$C'\to C$ such that $S$ has an analytic $\ell$-th root on the entire
inverse image of $A$.  A single such cover can be chosen for any finite
collection of node annuli.
\end{lemma}

\begin{proof}
Our aim is to make the expansion factors, the absolute values of the
integral slopes along edges, divisible by $\ell$ above the chosen edge.
On an inverse annulus this makes the Laurent exponent of
the pulled-back coordinate divisible by $\ell$; its constant factor and
the factor $1+h$ also have $\ell$-th roots by
Lemma~\ref{lem:annular-unit}.  We construct such a cover directly when
the edge lies on a circuit.  For a separating edge, we first pass to an
unramified cover in which every edge above it lies on a circuit.

We use the lifting theorem for tame coverings of metrized complexes
\cite[Theorem~7.7]{ABBR}.  A covering to which it applies consists of
finite residue-curve maps and a finite harmonic graph map: every fibre of
every marked point is represented by edge ends, their ramification
indices equal the edge expansion factors, and the component degree is
the sum of these factors over each base direction.  The residue maps
must be tamely ramified, and the augmented graph ramification must
vanish away from punctures
\cite[Definitions~2.19 and~2.21]{ABBR}.
We shall use no punctures.  The lift is then finite \'etale, and has
exactly the prescribed inverse images of the vertex set and skeleton
\cite[Definitions~4.25 and~4.31]{ABBR}.

First suppose that the edge $e$ corresponding to $A$ is nonseparating.
Choose a simple circuit containing $e$.  At each vertex on this circuit,
let $P,Q$ be the two circuit marks on its residue curve $D$.
Divisibility of $\operatorname{Pic}^0(D)$
\cite[Theorem~8.2]{MilneAV} supplies a degree-zero divisor
$E$ such that $P-Q+\ell E$ is principal.  Extract an $\ell$-th root of a
function with this divisor and take the smooth normalization.
The resulting cyclic cover of $D$ has degree $\ell$, is connected
because the order at $P$ is prime to $\ell$, and is totally ramified at
$P,Q$ and nowhere else.  Each other marked point has $\ell$ distinct
preimages.  At every vertex outside the circuit take $\ell$ identity
copies of its residue curve.

Above each circuit edge use one edge of expansion factor $\ell$ and of
length equal to the original length divided by $\ell$.  Above each
remaining edge use $\ell$ factor-one edges, matching all the unramified
marked fibres at its ends.  These rules also handle loops and chords
between circuit vertices.  All lengths belong to $\mathbb Q$.
At a lifted circuit vertex each base direction has total degree $\ell$;
at an identity-copy vertex it has total degree one.
The local Riemann--Hurwitz formula is
\[
 2g(D')-2=\ell\bigl(2g(D)-2\bigr)+2(\ell-1).
\]
The two extra terms are precisely those recorded by the two circuit
edges.  Thus the augmented graph ramification is zero at every vertex;
the identity copies contribute none.  All ramification indices are
prime to $p$.  The resulting tame covering is connected, because each
vertex outside the circuit can be joined by a lifted path to a circuit
vertex.

Lift these data by the cited theorem.  The inverse image of $A$ is
exactly the disjoint union of the node annuli above $e$.
For completeness, the inverse-skeleton equality excludes off-skeleton
points mapping to the skeleton.  Each off-skeleton open ball therefore
maps into a single off-skeleton ball downstairs, and continuity
identifies the images of their attachment points.  Retractions commute
on these balls as well as on the skeleton.  Since a node annulus is the
inverse image of its open edge under retraction, this proves the asserted
equality for the whole annulus.  On each annulus above $A$, the
pullback of $S$ has exponent $\ell$ or $-\ell$ in
Lemma~\ref{lem:annular-unit}.  It therefore has an analytic $\ell$-th
root on each of these disjoint annuli.

Now suppose that $e$ is separating.  Each side of the graph obtained by
cutting its interior contains a circuit or a positive-genus vertex.
Otherwise that side would be a genus-zero tree with, say, $N$ vertices.
Its total valence, including the bridge half-edge, would be $2N-1$,
contrary to stability, which requires at least $3N$.

Construct a connected unramified degree-$\ell$ covering of each side,
retaining its boundary mark for the eventual gluing.  If the side
contains a circuit, take a connected cyclic graph covering of degree
$\ell$, with identity residue-curve maps.  Such a graph covering is
obtained by a homomorphism of its graph fundamental group onto
$\mathbb Z/\ell\mathbb Z$.  If instead a positive-genus vertex is used,
take a connected cyclic \'etale cover of degree $\ell$ of its residue
curve, provided by a nonzero element of its Jacobian's $\ell$-torsion
\cite[Theorem~8.2 and Remark~8.4]{MilneAV}.  Explicitly, for a line
bundle $L$ of exact order $\ell$, a trivialization
$L^{\otimes\ell}\simeq\mathcal O$ makes
$\bigoplus_{j=0}^{\ell-1}L^{\otimes j}$ an algebra whose relative
spectrum is a cyclic \'etale cover: $\ell$ is invertible, and the exact
order of $L$ makes the cover connected.  Use $\ell$ identity copies at
the other vertices and match all marked fibres.  Paths to the
distinguished vertex show connectedness of the covered side.

Each covered side has $\ell$ preimages of the bridge mark.  Match them
in pairs to obtain $\ell$ distinct bridge edges, all with factor one.
Every such edge is now nonseparating: after removing one of them, the
two connected covered sides are still joined by the others.
The resulting covering of the original metrized complex is tame.
Indeed the component covers are unramified, all factors are one, and
Riemann--Hurwitz again gives zero graph ramification.  Lift it to a
finite \'etale curve cover.  Its triangulation remains stable: at a
vertex of local degree $d$, the quantity
$2g'-2+\operatorname{val}(v')$ is
$d(2g-2+\operatorname{val}(v))>0$.

Apply the nonseparating construction to each of these $\ell$ edges.
A connected component of the fibre product of the resulting finite
\'etale covers dominates each factor, so all the constructed roots
persist on a common cover.  On each factor-one annulus the pullback of
the original $S$ is a unit; Lemma~\ref{lem:annular-unit} supplies its root
once the new annular coordinate has a root.  This proves the assertion
also for separating edges.  Finally, the same finite fibre-product
argument treats finitely many original annuli simultaneously.
\end{proof}

\subsection{Moving disk collars into the stable skeleton}

The stable-edge lemma does not apply directly to a small disk around a
branch point: nodes introduced there by blowups need not survive
stabilization.  Resolution of nonsingularities supplies an \'etale
cover whose stable model maps to the chosen model with those nodes.
The next argument requires only that model morphism.

\begin{lemma}[\'Etale covers with roots on collars]\label{lem:etale-collars}
Let $C/\bar k$ be a smooth proper connected curve of genus at least two.
Let $b_1,\ldots,b_s\in C(\bar k)$ be distinct, and let $B_i$ be pairwise
disjoint sufficiently small coordinate disks in $C_\Omega^{\mathrm{an}}$
about $b_i$, with coordinates $t_i$ defined over $\bar k$ and vanishing
simply at $b_i$.  There is a connected finite \'etale cover
$Y\to C_\Omega$ and concentric open collars
\[
 A_i'=\{r_i<|t_i|<R_i\}\subset B_i,\qquad
 r_i,R_i\in|\Omega^\times|,\quad 0<r_i<R_i,
\]
such that $t_i$ has an analytic $\ell$-th root on the entire inverse
image of $A_i'$ in $Y^{\mathrm{an}}$.
\end{lemma}

\begin{proof}
All algebraic data descend to a finite extension of $k$.
After a further finite extension, choose a split semistable model
with the $b_i$ as disjoint smooth sections.  Shrinking the coordinate
disks causes no problem, since the desired collars need only lie inside
the originally specified disks.  In an integral local parameter $u_i$
at a section,
successive blowups at its reductions introduce node annuli
\[
 |\pi|^{j+1}<|u_i|<|\pi|^j .
\]
On a sufficiently small disk the quotient $t_i/u_i$ is an analytic
unit, since both coordinates have a simple zero at $b_i$.
Its absolute value is constant by Lemma~\ref{lem:annular-unit}.
Thus these annuli are also concentric annuli for $t_i$, and choosing
$j$ large puts them inside $B_i$.
Perform these modifications for all the sections and denote the
resulting proper semistable model by $\mathcal C$.
It need not be stable; the markings can now be forgotten.

The RNS theorem, applied with the set of all primes, and its
model-domination consequence
\cite[Theorem~A and Proposition~2.4(iv)]{MochizukiTsujimura}
give, after finite extension of the ground field, a finite \'etale
cover whose stable model maps to $\mathcal C$.
Pass to split reduction and a geometric connected component, and write
\[
 f:Y_0\longrightarrow C_\Omega,\qquad
 \mathcal Y_0\longrightarrow\mathcal C
\]
for the analytic-base-changed curve cover and the model morphism.
No finiteness assertion about the model morphism is needed.

Fix one of the target node annuli $A_i\subset B_i$ just constructed,
and choose an irrational real number $a$ strictly between its two
additive endpoint valuations.  Consider the entire inverse image
\[
 f^{-1}\bigl(\{x\in A_i:-\log_p|t_i(x)|=a\}\bigr).
\]  Reduction is compatible with the model
morphism.  A point reducing to a component generic point of
$\mathcal Y_0$ is a model vertex, hence is type~2 and gives $t_i$ a
value in $|\Omega^\times|$; its additive valuation is rational.
It cannot lie above this level.  A point reducing to a smooth closed
point lies in an open disk tube.  If it lies above the chosen level,
that smooth point maps to the target node, so the entire disk tube
maps into $A_i$.  The pullback of $t_i$ is a unit there.
By Lemma~\ref{lem:annular-unit} its norm is constant and belongs to
$|\Omega^\times|$, which again excludes the irrational level.
All remaining points above the level lie in source node annuli whose
nodes map to the target node.

Let $N_i$ be the union of all source node annuli whose nodes map to
the target node defining $A_i$; this is a finite union.
The complement of $N_i$ in the proper analytic curve $Y_0^{\mathrm{an}}$
is compact.  Intersect it with the inverse image of a closed
concentric subannulus of $A_i$ containing the chosen level in its
interior.  On this compact intersection the continuous real function
$-\log_p|t_i|$ has compact image avoiding $a$.
Consequently there is an open interval about $a$ with rational
endpoints whose entire inverse image lies in $N_i$.
It determines a collar $A_i'\subset A_i$.

Every source node annulus in $N_i$ maps into $A_i$ by reduction
compatibility, so the pullback of $t_i$ is an analytic unit on it.
Apply Lemma~\ref{lem:stable-edge-root} simultaneously to all these
source annuli, for all $i$.  Their coordinates acquire $\ell$-th roots
on a common connected finite \'etale cover of $Y_0$.
Lemma~\ref{lem:annular-unit} then gives roots of the pullbacks of
the $t_i$ on the full inverse images of the $A_i'$.
The resulting composite cover of $C_\Omega$ is finite \'etale.
\end{proof}

\begin{proposition}[The first stage]\label{prop:first-stage}
The field $L_\ell$ is contained in $H$.
\end{proposition}

\begin{proof}
Let $F/K_0$ be a finite Galois extension whose horizontal ramification
indices divide $\ell$, and let $C\to X_{\bar k}$ be its smooth
proper curve cover.  If it is unramified, its function field is already
contained in $\widetilde K$.  Otherwise take disjoint branch-point
disks avoiding $\xi$, small enough for
Lemma~\ref{lem:local-normal-form}.
Lemma~\ref{lem:etale-collars}, applied on $X_{\bar k}$, supplies a finite
\'etale $Y\to X_\Omega$ and collars on whose full inverse images the
branch coordinates have $\ell$-th roots.

Apply Proposition~\ref{prop:patching} with these roots.  It gives a
connected finite \'etale cover $V_\Omega\to Y$ and an analytic map to
$C_\Omega$ over the exterior of the smaller branch disks.
The composite $V_\Omega\to X_\Omega$ is finite \'etale.
Proper base-extension invariance of finite \'etale covers descends it,
up to isomorphism over $X_\Omega$, to a cover $V\to X_{\bar k}$.
Here we use \cite[Expos\'e~X, Corollary~1.8]{SGA1}.
Its function field has a $K_0$-embedding into
$\widetilde K\subset H$.
Lemma~\ref{lem:point-transfer} therefore yields a $K_0$-embedding
$F\hookrightarrow H$ at the fixed $H$-valued point above $\xi$.
Since $F/K_0$ is normal and both fields lie in $E$, this says
$F\subset H$ for the chosen copy of $F$.
Taking their compositum proves the proposition.
\end{proof}
